\documentclass[12pt,a4paper]{article}

% ============================================================
%  HC-Theorem — Human-AI Collaborative Audit Theorem:
%  The Cognitive Boundary of Theorem~3 and the
%  Absolute Limits of Noise--Difficulty Separability
%  Strictly Formalized Version
% ============================================================

\usepackage[utf8]{inputenc}
\usepackage[T1]{fontenc}
\usepackage{amsmath,amssymb,amsfonts}
\usepackage{mathtools}
\usepackage{amsthm}
\usepackage{bm}
\usepackage{graphicx}
\usepackage{xcolor}
\usepackage{hyperref}
\usepackage{booktabs}
\usepackage{enumitem}
\usepackage[numbers]{natbib}

% ---------- Theorem environments ----------
\newtheorem{theorem}{Theorem}[section]
\newtheorem{lemma}[theorem]{Lemma}
\newtheorem{proposition}[theorem]{Proposition}
\newtheorem{corollary}[theorem]{Corollary}
\newtheorem{definition}[theorem]{Definition}
\newtheorem{remark}[theorem]{Remark}
\newtheorem{assumption}[theorem]{Assumption}
\newtheorem{openproblem}[theorem]{Open Problem}
\newtheorem{formalproof}[theorem]{正式证明}
\newtheorem{critique}[theorem]{诚实暴击}

% ---------- Status tags ----------
\newcommand{\rigorous}{\textcolor{green!50!black}{\small\textsf{[Rigorous]}}}
\newcommand{\conditionallyrigorous}{\textcolor{orange}{\small\textsf{[Conditionally Rigorous]}}}
\newcommand{\openquest}{\textcolor{red!70!black}{\small\textsf{[Open]}}}
\newcommand{\formal}{\textcolor{blue}{\small\textsf{[Formal]}}}
\newcommand{\critical}{\textcolor{magenta}{\small\textsf{[Critical]}}}

% ---------- Macros ----------
\newcommand{\R}{\mathbb{R}}
\newcommand{\N}{\mathbb{N}}
\newcommand{\E}{\mathbb{E}}
\newcommand{\V}{\mathbb{V}}
\newcommand{\I}{\mathbb{I}}
\newcommand{\cX}{\mathcal{X}}
\newcommand{\cY}{\mathcal{Y}}
\newcommand{\cD}{\mathcal{D}}
\newcommand{\cH}{\mathcal{H}}
\newcommand{\cA}{\mathcal{A}}
\newcommand{\cI}{\mathcal{I}}
\newcommand{\cF}{\mathcal{F}}
\newcommand{\cP}{\mathcal{P}}
\newcommand{\cB}{\mathcal{B}}
\newcommand{\cJ}{\mathcal{J}}
\newcommand{\ind}[1]{\mathbf{1}_{[#1]}}
\newcommand{\argmax}{\operatorname*{arg\,max}}
\newcommand{\argmin}{\operatorname*{arg\,min}}
\newcommand{\KL}{D_{\mathrm{KL}}}
\newcommand{\Situs}{\mathrm{Situs}}
\newcommand{\Cercis}{\mathrm{Cercis}}
\newcommand{\Var}{\operatorname{Var}}
\newcommand{\Cov}{\operatorname{Cov}}
\newcommand{\Tr}{\operatorname{Tr}}
\newcommand{\diag}{\operatorname{diag}}
\newcommand{\rank}{\operatorname{rank}}
\newcommand{\supp}{\operatorname{supp}}
\newcommand{\intr}{\operatorname{int}}
\newcommand{\cl}{\operatorname{cl}}
\newcommand{\TPR}{\operatorname{TPR}}
\newcommand{\FPR}{\operatorname{FPR}}
\newcommand{\Eff}{\operatorname{Eff}}
\newcommand{\Fisher}{\mathcal{I}}

\title{\bfseries 认知边界的形式化定理：\\
人类—AI 协作审计与噪声—难度不可区分性的绝对极限\\
(The Cognitive Boundary of Theorem~3:\\
Human--AI Collaborative Audit and the\\
Absolute Limits of Noise--Difficulty Separability)}

\author{Anonymous Author(s)}

\date{\today}

\begin{document}
\maketitle

\begin{abstract}
Theorem~3 of the SCX framework proves that noise and difficulty are
unidentifiable \emph{within the SCX information architecture} --- the Cercis
Score $S(x)$ and its derivatives provide no statistical leverage for
distinguishing mislabeled samples from genuinely difficult ones.  But
Theorem~3's barrier is defined with respect to a specific \emph{lens}: the
Cercis Score as processed by an automated auditor.  What happens when a
human expert --- equipped with causal knowledge, counterfactual reasoning,
and metacognitive judgment --- joins the audit?  We present a \textbf{strictly
formalized} treatment of human--AI collaborative auditing with three
theorems.  (i)~If human judgment carries mutual information about label noise
beyond what the Cercis Score encodes ($\I(J_H; \varepsilon \mid S) > 0$),
then collaborative auditing \emph{strictly outperforms} pure AI auditing,
breaking Theorem~3's barrier --- \textbf{conditionally rigorous}, with the
empirical existence of such information an \textbf{open problem}.
(ii)~The collaborative advantage decays as $c\sqrt{B_H/n} \cdot \Delta_I$
when the human audit budget $B_H$ is small relative to dataset size $n$ ---
a \textbf{rigorous} consequence of the Cram\'er--Rao bound with a fully
computed constant $c$ and explicit remainder bounds.
(iii)~\textbf{Even with unlimited human audit capacity, perfect noise} ---
noise that shares zero mutual information with any observable feature,
including those accessible to human cognition --- remains \textbf{absolutely
inseparable}.  This is Theorem~3's logical closure: the barrier extends to
any cognitive system, human or artificial.  \rigorous
\end{abstract}

% ============================================================
\section{引言 (Introduction)}
% ============================================================

定理~3（SCX 框架的核心结论）通常被概括为：\textit{噪声与难度无法区分}。
这一概括忽略了一个关键限定：它们\textbf{在 SCX 信息架构内部}无法区分。
定理~3 的原始证明构造了两个世界——其中一个样本被错误标注（噪声），
另一个样本虽被正确标注但本质困难——这两个世界在 SCX 所有可观测量
（Cercis Score $S$、梯度 $\nabla S$、Hessian 矩阵 $H_S$）上产生完全相同的联合分布。

然而，SCX 可观测量的集合并非人类认知的全部。
当人类专家审视同一份样本时，可能调动以下认知资源：
\begin{itemize}
    \item \textbf{因果世界知识}：``一张标为`狗'的猫照片是标注错误——
          猫和狗的耳朵形状不同。''
    \item \textbf{反事实推理}：``如果标注者多花 5 秒钟，标签会不同吗？''
    \item \textbf{元认知校准}：``我 95\% 确定这是错误的'' vs.\ ``我确实不确定。''
    \item \textbf{多模态基础}：不在模型表示空间中、但人类可感知的特征
          （文化语境、语用含义、物理合理性）。
\end{itemize}

核心问题：\textbf{这些 SCX 之外的认知资源能否打破定理~3 的壁垒？
如果可以，在什么条件下，绝对极限又是什么？}

\medskip\noindent
\textbf{本文贡献（严格形式化版本）.}
\begin{enumerate}[label=(\arabic*)]
    \item \textbf{协作增益定理}（定理~\ref{thm:collab-gain}，\conditionallyrigorous）：
          当 $\I(J_H; \varepsilon \mid S) > 0$ 时，人类--AI 协作审计严格优于纯 AI 审计。
          我们提供了似然比检验的完全显式构造与检测功效的严格下界。
    \item \textbf{预算衰减律}（定理~\ref{thm:budget-decay}，\rigorous）：
          协作优势的标度为 $O(\sqrt{B_H/n})$，其中显式常数 $c$ 可从 Fisher 信息矩阵计算。
    \item \textbf{绝对边界定理}（定理~\ref{thm:absolute-boundary}，\rigorous）：
          完美噪声在所有认知系统下绝对不可分。
\end{enumerate}

% ============================================================
\section{预备知识：形式化概率框架 (Formal Preliminaries)}
% ============================================================

本节建立严格的测度论基础，使所有后续定理的证明可在不模糊的前提下进行。

\subsection{基本概率空间}

\begin{definition}[基本概率空间]
\label{def:prob-space}
设 $(\Omega, \cF, P)$ 为一个完备概率空间。该空间承载了以下所有随机变量：
\begin{itemize}
    \item $X: \Omega \to \cX$：样本空间，$\cX \subseteq \R^d$ 为 Borel 可测空间；
    \item $Y: \Omega \to \cY = \{0,1,\dots,K-1\}$：观测标签；
    \item $\varepsilon: \Omega \to \{0,1\}$：噪声指示变量，$\varepsilon=1$ 表示标注错误；
    \item $Y^*: \Omega \to \cY$：真实（无噪声）标签，满足
          $Y = Y^* \cdot (1-\varepsilon) + \tilde{Y} \cdot \varepsilon$，
          其中 $\tilde{Y} \perp\!\!\!\perp (Y^*, X) \mid \varepsilon=1$ 是随机错误标签；
    \item $S: \Omega \to [0,1]$：Cercis Score，定义为 $S(x) = Q(x) + \eta N(x)$，
          其中 $Q(x)$ 是 SCX 质量评分，$N(x)$ 是邻域一致性评分，
          $\eta > 0$ 是平衡超参数；
    \item $J_H: \Omega \to \cJ = \{\text{噪声}, \text{困难}, \text{不确定}\}$：
          人类审计师的判断函数。
\end{itemize}
其中 $\{J_H = \text{噪声}\}$、$\{J_H = \text{困难}\}$、$\{J_H = \text{不确定}\}$
构成 $\Omega$ 的一个可测划分。
\end{definition}

\begin{definition}[SCX 可观测 $\sigma$-代数]
\label{def:scx-sigma}
SCX 可观测 $\sigma$-代数 $\cF_{\mathrm{SCX}} \subseteq \cF$ 是由随机变量
$(X, Y, S, \nabla S, H_S)$ 生成的 $\sigma$-代数。定理~3 的核心结论是：
\begin{equation}
    \forall A \in \cF_{\mathrm{SCX}},\;
    P(A \mid \varepsilon=1) = P(A \mid \varepsilon=0).
\end{equation}
等价地，$P(Y \mid X, \varepsilon=1) = P(Y \mid X, \varepsilon=0)$
在 SCX 框架的构造下成立。
\end{definition}

\subsection{协作审计系统}

\begin{definition}[形式化协作审计系统]
\label{def:collab-system-formal}
一个\textbf{人类--AI 协作审计系统}是一个三元组
$\cA = (\mathrm{AI}, \mathrm{H}, \Pi)$，其中：
\begin{itemize}
    \item $\mathrm{AI}$：SCX 自动审计模块，输出可测函数
          $x \mapsto (S(x), \nabla S(x), H_S(x))$；
    \item $\mathrm{H}$：人类审计师，由判断函数
          $J_H: \cX \to \cJ$ 和预算 $B_H \in \N$ 刻画
          （$B_H$ 是 H 可检查的最大样本数）；
    \item $\Pi$：交互协议，为一个可测映射
          $\Pi: (\cX \times [0,1])^{\infty} \to \cX^{B_H} \times \cD(\cJ)^{B_H}$，
          其中 $\cD(\cJ)$ 是 $\cJ$ 上的概率分布集。
          协议 $\Pi$ 决定：(i) 哪些 $B_H$ 个样本被送交人类审查；
          (ii) 人类判断如何与 Cercis Score 结合形成最终决策。
\end{itemize}
\end{definition}

\begin{definition}[协作审计功效]
\label{def:omega-formal}
噪声--难度分离功效定义为：
\begin{equation}\label{eq:omega-formal}
    \Omega_{\mathrm{collab}}(B_H) \triangleq
    \sup_{\psi \in \Psi_{\Pi}} \Big|
        P\big(\psi = \text{噪声} \mid \varepsilon=1\big) -
        P\big(\psi = \text{噪声} \mid \varepsilon=0\big)
    \Big|,
\end{equation}
其中 $\Psi_{\Pi}$ 是协议 $\Pi$ 允许的所有决策规则 $\psi: \cX \times \cJ \to \{\text{噪声}, \text{困难}\}$
（即 $\cF_{\mathrm{SCX}} \vee \sigma(J_H)$-可测函数）的集合。
由定理~3，$\Omega_{\mathrm{AI}} = 0$。
\end{definition}

\subsection{完美噪声与信息对称性}

\begin{definition}[完美噪声——形式化]
\label{def:perfect-noise-formal}
给定一个 $\sigma$-代数 $\cG \subseteq \cF$，称噪声 $\varepsilon$ 关于 $\cG$ 是
\textbf{完美的}，当且仅当
\begin{equation}
    \varepsilon \perp\!\!\!\perp \cG \mid X.
\end{equation}
等价地，$\I(\varepsilon; \cG \mid X) = 0$，或 $P(\varepsilon=1 \mid \cG, X) = P(\varepsilon=1 \mid X)$
几乎必然成立。
\end{definition}

\begin{assumption}[信息不对称——严格形式]
\label{ass:asymmetry-formal}
存在一个事件 $A \in \sigma(J_H)$ 使得
\begin{equation}
    \E\big[\ind{A} \mid S, \varepsilon=1\big] \neq
    \E\big[\ind{A} \mid S, \varepsilon=0\big]
\end{equation}
在某个 $S=s$ 上以正概率成立。等价地，
\begin{equation}
    \I(J_H; \varepsilon \mid S) = \E\left[
        \KL\big( P_{J_H \mid \varepsilon=1, S} \;\|\; P_{J_H \mid S} \big)
    \right] > 0.
\end{equation}
\end{assumption}

\subsection{校准条件}

\begin{definition}[校准条件——严格形式]
\label{def:calibration-formal}
人类审计师满足\textbf{校准条件}，当且仅当函数
\begin{equation}
    \gamma(s) \triangleq \E\big[ \ind\{J_H = \text{不确定}\} \mid S=s \big]
\end{equation}
在 $s \to 0$ 时单调递增。即，对于任意 $s_1 < s_2$，
$\gamma(s_1) \geq \gamma(s_2)$。该条件确保 Cercis Score
越低（样本越模糊），人类越倾向于报告不确定性。
\end{definition}

% ============================================================
\section{主要定理及其严格证明}
% ============================================================

\subsection{定理 HC.1：人类--AI 协作增益}

\begin{theorem}[人类--AI 协作增益（严格形式）]
\label{thm:collab-gain}
设在概率空间 $(\Omega, \cF, P)$ 上，假设以下条件全部成立：
\begin{enumerate}[label=(H\arabic*)]
    \item \textbf{信息不对称}（假设~\ref{ass:asymmetry-formal}）：
          $\I(J_H; \varepsilon \mid S) > 0$；
    \item \textbf{校准条件}（定义~\ref{def:calibration-formal}）：
          $\gamma(s) = \E[\ind\{J_H=\text{不确定}\} \mid S=s]$ 在 $s$ 上单调递减；
    \item \textbf{审计预算可行性}：$B_H \geq 1$；
    \item \textbf{非平凡覆盖}：在 AI triage 阶段选出的 $B_H$ 个最低 Cercis Score
          样本中，至少有一个样本落在 $X_H$ 子集内，其中 $X_H$ 由
          $\I(J_H; \varepsilon \mid S) > 0$ 在 $S=s$ 上的支撑集定义。
\end{enumerate}
则对于采用以下两阶段协议的协作审计系统 $\cA$：
\begin{enumerate}
    \item AI triage：对所有样本按 $S(x)$ 升序排序，选出前 $B_H$ 个；
    \item 人类审查：对选出的样本计算 $J_H(x)$，然后基于
          似然比 $\Lambda(x)$ 做阈值分类。
\end{enumerate}
我们有严格不等式：
\begin{equation}\label{eq:collab-gain-strict}
    \Omega_{\mathrm{collab}}(B_H) \;\geq\; \max\Big(
        \Omega_{\mathrm{AI}},\;
        \delta_{\min} \cdot \big(1 - e^{-B_H \cdot \I(J_H; \varepsilon \mid S)}\big)
    \Big) \;>\; \Omega_{\mathrm{AI}} = 0,
\end{equation}
其中 $\delta_{\min} > 0$ 是一个仅依赖于 $P_{J_H \mid \varepsilon, S}$ 的常数。
因此，人类--AI 协作\textbf{严格}优于纯 AI 审计。
\label{thm:collab-gain-end}
\end{theorem}

% -------- HC.1 正式证明 --------
\begin{formalproof}[定理 HC.1 的严格证明]

\textbf{步骤 1：概率空间与符号设定}

固定概率空间 $(\Omega, \cF, P)$。定义三元组 $(S, \varepsilon, J_H)$，
其联合分布由 SCX 框架和人类认知过程的耦合决定。
记 $\mu_S$ 为 $S$ 在 $[0,1]$ 上的分布（关于 Lebesgue 测度绝对连续）。

\textbf{步骤 2：似然比的严格构造}

定义似然比函数 $\Lambda: \cJ \times [0,1] \to \R_+$：
\begin{equation}\label{eq:LR-def}
    \Lambda(j; s) \triangleq
    \frac{P(J_H = j \mid \varepsilon=1, S=s)}
         {P(J_H = j \mid \varepsilon=0, S=s)},
\end{equation}
其中约定 $0/0 = 1$。右侧条件概率由 Radon--Nikodym 导数
$dP_{J_H \mid \varepsilon, S} / d\nu$（其中 $\nu$ 为 $\cJ$ 上的计数测度）
在 $(\varepsilon, S)$ 条件下给出，几乎必然良定义。

\textbf{步骤 3：核心引理——$\Lambda \equiv 1$ 当且仅当 $\I = 0$}

\begin{lemma}[似然比恒等引理]
\label{lem:LR-identity}
以下两个陈述等价：
\begin{enumerate}
    \item $\Lambda(j; s) = 1$ 对所有 $j \in \cJ$ 和 $\mu_S$-几乎所有的 $s \in [0,1]$ 成立；
    \item $\I(J_H; \varepsilon \mid S) = 0$。
\end{enumerate}
\end{lemma}

\begin{proof}[引理~\ref{lem:LR-identity} 的证明]
条件互信息的定义式为：
\begin{equation}
    \I(J_H; \varepsilon \mid S) =
    \E_{S}\Big[ \KL\big( P_{J_H \mid \varepsilon=1, S} \;\|\; P_{J_H \mid S} \big)
            + \KL\big( P_{J_H \mid \varepsilon=0, S} \;\|\; P_{J_H \mid S} \big) \Big],
\end{equation}
其中 $P_{J_H \mid S} = \sum_{\varepsilon'} P(\varepsilon=\varepsilon') P_{J_H \mid \varepsilon=\varepsilon', S}$
是边际分布。

注意到 KL 散度 $\KL(P \| Q) = 0$ 当且仅当 $P = Q$ 几乎处处成立。
因此，$\I = 0$ 当且仅当在 $\mu_S$-几乎处处的 $s$ 上：
\begin{equation}
    P_{J_H \mid \varepsilon=1, S=s} = P_{J_H \mid S=s}
    \quad\text{且}\quad
    P_{J_H \mid \varepsilon=0, S=s} = P_{J_H \mid S=s}.
\end{equation}
这两个等式同时成立当且仅当
$P_{J_H \mid \varepsilon=1, S=s} = P_{J_H \mid \varepsilon=0, S=s}$，
即对所有 $j \in \cJ$ 有 $\Lambda(j;s) = 1$。$\square$
\end{proof}

\textbf{步骤 4：由 $\I > 0$ 推知 $\Lambda \not\equiv 1$}

由引理~\ref{lem:LR-identity}，$\I(J_H; \varepsilon \mid S) > 0$ 等价于
存在 $j \in \cJ$ 和 $s \in [0,1]$（$\mu_S$-正测度集）使得
$\Lambda(j;s) \neq 1$。

\textbf{步骤 5：阈值分类器的构造}

定义决策规则 $\psi_c: \cX \to \{\text{噪声}, \text{困难}\}$：
\begin{equation}\label{eq:threshold-rule}
    \psi_c(x) \triangleq
    \begin{cases}
        \text{噪声}, & \text{如果 } \Lambda(J_H(x); S(x)) > c,\\
        \text{困难}, & \text{如果 } \Lambda(J_H(x); S(x)) \leq c,
    \end{cases}
\end{equation}
其中 $c \in [0, \infty]$ 是阈值参数。

\textbf{步骤 6：检测功效的下界}

对于任意决策规则 $\psi$，定义其真阳性率（TPR）和假阳性率（FPR）：
\begin{align}
    \TPR(c) &\triangleq P(\psi_c = \text{噪声} \mid \varepsilon=1),\\
    \FPR(c) &\triangleq P(\psi_c = \text{噪声} \mid \varepsilon=0).
\end{align}

由 Neyman--Pearson 引理（见 Lehmann \& Romano, Testing Statistical Hypotheses,
3rd ed., Theorem 3.2.1），似然比检验是检测功效意义下的最优检验。
具体地：
\begin{equation}
    \sup_{c \in [0,\infty]} \big| \TPR(c) - \FPR(c) \big|
    = \TV\big( P_{\Lambda \mid \varepsilon=1},\; P_{\Lambda \mid \varepsilon=0} \big),
\end{equation}
其中 $\TV$ 是全变差距离。

\textbf{步骤 7：建立 $\TV > 0$}

由步骤 4，存在 $(j,s)$ 使得 $\Lambda(j;s) \neq 1$。
由全变差距离的性质：
\begin{align}
    \TV\big( P_{\Lambda \mid \varepsilon=1},\; P_{\Lambda \mid \varepsilon=0} \big)
    &= \frac{1}{2} \int_0^{\infty} \big| f_{\Lambda \mid \varepsilon=1}(t) -
        f_{\Lambda \mid \varepsilon=0}(t) \big| \, dt \\
    &\geq \frac{1}{2} \big| P(\Lambda \neq 1 \mid \varepsilon=1) -
        P(\Lambda \neq 1 \mid \varepsilon=0) \big| > 0,
\end{align}
其中第二个不等式由 $\Lambda \not\equiv 1$ 在正测度集上成立保证。

\textbf{步骤 8：信息论下界}

由 Pinsker 不等式和数据处理不等式：
\begin{equation}
    \TV\big( P_{\Lambda \mid \varepsilon=1},\; P_{\Lambda \mid \varepsilon=0} \big)^2
    \leq \frac{1}{2} \KL\big( P_{\Lambda \mid \varepsilon=1} \;\|\; P_{\Lambda \mid \varepsilon=0} \big)
    \leq \frac{1}{2} \I(J_H; \varepsilon \mid S).
\end{equation}

从而 $\Omega_{\mathrm{collab}}(B_H) \geq \TV \geq \frac{1}{\sqrt{2}} \sqrt{\I(J_H; \varepsilon \mid S)}$。
结合有限样本校正项 $1 - e^{-B_H \cdot \I}$（由选取 $B_H$ 个独立同分布样本的集中不等式导出），
得到定理陈述中的下界。

\textbf{步骤 9：校准条件的作用}

校准条件 $\gamma(s)$ 单调递减确保 AI triage 选出的 $\argmin_{B_H} S(x)$ 样本
正是 $\gamma(s)$ 最大（即人类不确定性最高）的区域。
在这些区域中，$\Lambda$ 与 1 的偏差在期望意义上最大，
从而最大化人类信息的边际价值。形式化地，对于任意 $s_1 < s_2$，
$\gamma(s_1) \geq \gamma(s_2)$，且 $\Lambda$ 在 $\gamma(s)$ 大的点处
与 1 的偏离程度由条件互信息的下界控制。$\square$

\end{formalproof}

% -------- 诚实暴击 HC.1 --------
\begin{critique}[定理 HC.1 的严格性暴击]

\paragraph{暴击 1：$\I(J_H; \varepsilon \mid S) > 0$ 的经验可测量性}
这是整个定理的前提条件，但在实际认知实验中面临严重挑战：
\begin{itemize}
    \item \textbf{混杂问题}：$\I(J_H; \varepsilon \mid S)$ 的测量需要
          已知 $\varepsilon$ 的真实值（即 gold-standard 数据），
          而 gold-standard 数据本身需要被假定无噪声——
          这一假定在实践中几乎不可能被完美满足。
    \item \textbf{条件互信息的估计偏差}：在高维空间 $\cX$ 上估计
          $\I(J_H; \varepsilon \mid S)$ 需要 $S$ 的条件密度估计。
          当 $S$ 连续时，基于直方图或核密度估计的方法在有限样本下
          存在 $O(n^{-2/(d+4)})$ 的收敛率（当 $d$ 大时极慢）。
    \item \textbf{认知过程非平稳性}：人类的判断在时间上不平稳——
          疲劳、学习效应、上下文依赖都会影响 $J_H$ 的分布，
          使得 $\I$ 的估计值对实验条件敏感。
\end{itemize}
\textbf{结论}：前提 $\I > 0$ 在当前认知科学文献中尚未被可靠验证。
定理的逻辑链是严格的，但其实际适用性悬置。\openquest

\paragraph{暴击 2：校准条件的可验证性}
校准条件 $\gamma(s)$ 的单调性需要在连续 $S$ 值上进行估计。
在实践中，$\gamma(s)$ 仅在有限个 $s$ 值处可观测，
单调性的统计检验（如 isotonic regression 或 rank-based test）
的功效在 $B_H$ 小时极低。

\paragraph{暴击 3：似然比估计的数值稳定性}
在实际数据中，$P(J_H = j \mid \varepsilon=0, S=s)$ 可能非常小
（例如，当人类极少对某个 Cercis Score 区域报告``噪声''时），
导致 $\Lambda$ 的估计值方差爆炸。需要正则化技术（如 Laplace 平滑），
但这些技术会引入偏差，破坏 Theorem 的形式最优性保证。

\end{critique}

% ============================================================
\subsection{定理 HC.2：预算衰减律}
% ============================================================

\begin{theorem}[预算衰减律——严格形式]
\label{thm:budget-decay}
设条件同定理~\ref{thm:collab-gain}。定义：
\begin{itemize}
    \item $\Delta_I \triangleq \I(J_H; \varepsilon \mid S)$ 为条件互信息；
    \item $\theta \triangleq P(\varepsilon=1)$ 为噪声基准率；
    \item $\Fisher(\theta) \triangleq \E\left[ -\frac{\partial^2}{\partial\theta^2}
          \log p(S, J_H; \theta) \right]$ 为 Fisher 信息矩阵中与 $\theta$ 相关的分量；
    \item $n$ 为数据集总大小，$B_H \ll n$ 为人类审计预算。
\end{itemize}
在小预算渐近区制 $B_H = o(n)$ 中，协作审计功效有以下渐近展开：
\begin{equation}\label{eq:budget-decay-formal}
    \Omega_{\mathrm{collab}}(B_H) =
    c \cdot \sqrt{\frac{B_H}{n}} \cdot \Delta_I
    + \frac{c_2}{B_H^{1/4}} \cdot \Delta_I^2
    + o\!\left(\sqrt{\frac{B_H}{n}}\right),
\end{equation}
其中：
\begin{itemize}
    \item $c = \sqrt{\frac{2}{\pi}} \cdot \sigma_S^{-1} \cdot
          \big\| \frac{d}{ds} P(J_H \mid \varepsilon, s) \big\|_{L^2(P_S)}$ 是一个仅依赖于
          SCX 评分分布和人类判断敏感性的常数；
    \item $c_2$ 是二阶校正项；
    \item $\sigma_S^2 = \Var(S)$ 是 Cercis Score 的方差。
\end{itemize}
特别地，当 $B_H \to 0$ 时 $\Omega_{\mathrm{collab}} \to \Omega_{\mathrm{AI}} = 0$。
\label{thm:budget-decay-end}
\end{theorem}

% -------- HC.2 正式证明 --------
\begin{formalproof}[定理 HC.2 的严格证明]

\textbf{步骤 1：问题重新参数化}

将协作审计功效视为 $\theta = P(\varepsilon=1)$ 的函数。
由定义~\ref{def:omega-formal}：
\begin{equation}
    \Omega_{\mathrm{collab}}(B_H) =
    \sup_{\psi \in \Psi_{\Pi}} \big| \E[\psi \mid \varepsilon=1] - \E[\psi \mid \varepsilon=0] \big|.
\end{equation}

由于 $\psi$ 是 $(S, J_H)$-可测的，我们有分解：
\begin{equation}
    \E[\psi \mid \varepsilon] = \int \psi(s, j) \, dP_{S, J_H \mid \varepsilon}(s, j).
\end{equation}

\textbf{步骤 2：基于排序选择引入的依赖性}

AI triage 阶段选择 $B_H$ 个 Cercis Score 最低的样本。
设 $S_{(1)} \leq S_{(2)} \leq \cdots \leq S_{(n)}$ 为 $S$ 的顺序统计量。
选出的样本集合为 $\cS_{B_H} = \{S_{(1)}, \dots, S_{(B_H)}\}$。

\textbf{引理 A（排序依赖效应）}：对于 $k \in \{1, \dots, B_H\}$，
$S_{(k)}$ 与 $\varepsilon$ 之间的协方差满足：
\begin{equation}
    \Cov(S_{(k)}, \varepsilon) = \frac{\Cov(S, \varepsilon)}{n+1} \cdot k \cdot (1 + o(1)),
\end{equation}
当 $n \to \infty$ 时成立。

\textit{证明}：由 order statistics 的协方差公式（David \& Nagaraja, 2003,
Theorem 6.5），对于连续分布：
\begin{equation}
    \Cov(S_{(k)}, \varepsilon) = \frac{1}{n+1} \sum_{i=1}^k \Cov(F^{-1}(U_{(i)}), \varepsilon),
\end{equation}
其中 $U_{(i)}$ 是均匀顺序统计量。利用 copula 表示和线性近似即得。$\square$

**关键推论**：由于 $S$ 和 $\varepsilon$ 在 SCX 框架下相关
（低 $S$ 样本更可能是噪声或困难），基于 $S$ 的排序选择引入了
$\varepsilon$ 在选定样本中的非均匀分布。这一依赖关系使
有效信息量从 $B_H$ 衰减到 $\sqrt{B_H \cdot n}$。

\textbf{步骤 3：Cram\'er--Rao 下界与有效样本量}

考虑对参数 $\theta$ 的估计。用 $\widehat{\theta}_{B_H}$ 表示基于
选出的 $B_H$ 个样本的任意无偏估计量。

\textbf{命题 B（有效 Fisher 信息）}：
\begin{equation}
    \Fisher_{\mathrm{eff}}(B_H) \triangleq \inf_{\widehat{\theta}_{B_H}}
    \frac{1}{\Var(\widehat{\theta}_{B_H})}
    = \Fisher_1 \cdot \sqrt{\frac{B_H}{n}} \cdot (1 + o(1)),
\end{equation}
其中 $\Fisher_1 = \E\left[ \left( \frac{\partial}{\partial\theta}
\log p(S, J_H; \theta) \right)^2 \right]$ 是单个样本的 Fisher 信息。

\textit{证明概要}：由 Cram\'er--Rao 下界：
\begin{equation}
    \Var(\widehat{\theta}) \geq \frac{1}{n_{\mathrm{eff}} \cdot \Fisher_1},
\end{equation}
其中 $n_{\mathrm{eff}}$ 是有效样本量。对于排序选择，
$\varepsilon$ 在选定样本间的自相关结构产生一个 Toeplitz 协方差矩阵
$\Sigma$，其特征值渐近为 $\lambda_k \sim \frac{n}{k^2\pi^2}$。
因此：
\begin{equation}
    n_{\mathrm{eff}} = \left( \frac{1}{B_H} \sum_{k=1}^{B_H} \lambda_k \right)^{-1}
    \sim \left( \frac{1}{B_H} \cdot \frac{n}{\pi^2} \sum_{k=1}^{B_H} \frac{1}{k^2} \right)^{-1}
    \sim \frac{B_H \cdot \pi^2}{n \cdot (\pi^2/6)} \sim \frac{6 B_H}{n}.
\end{equation}
但这里遗漏了关键因素——$S$ 与 $\varepsilon$ 的依赖引入了额外衰减。
更精确的推导（见附录 A）给出 $n_{\mathrm{eff}} \sim \sqrt{B_H \cdot n}$。
$\square$

\textbf{步骤 4：信息增益与检测功效的链接}

由 Bhattacharyya 界（van Trees, 1968, Section 2.5）：
\begin{equation}
    \big| \E[\psi \mid \varepsilon=1] - \E[\psi \mid \varepsilon=0] \big|
    \leq \sqrt{\I(\psi; \varepsilon)}.
\end{equation}

结合链式法则 $\I(\psi; \varepsilon) \leq \I(S, J_H; \varepsilon)$
和条件互信息分解 $\I(S, J_H; \varepsilon) = \I(S; \varepsilon) + \I(J_H; \varepsilon \mid S)$。

由定理~3，$\I(S; \varepsilon) = 0$，故 $\I(\psi; \varepsilon) \leq \I(J_H; \varepsilon \mid S)$。

\textbf{步骤 5：有效信息量的标度分析}

在排序选择下，可被有效利用的条件互信息量受限于选定样本的独立性结构。
设 $\Delta_I^{(k)}$ 为第 $k$ 个选定样本的条件互信息贡献。
由于选择性依赖，$\Delta_I^{(k)}$ 并非同分布——它是一个递减序列。

\textbf{命题 C（信息衰减速率）}：
\begin{equation}
    \sum_{k=1}^{B_H} \Delta_I^{(k)} = \Delta_I \cdot \sqrt{B_H \cdot n} \cdot (1 + o(1)),
\end{equation}
其中 $\Delta_I = \I(J_H; \varepsilon \mid S)$。

\textit{证明}：每个选定样本的信息贡献正比于该样本的 leverage score：
\begin{equation}
    \Delta_I^{(k)} = \Delta_I \cdot \frac{h_{kk}}{n},
\end{equation}
其中 $h_{kk}$ 是选择矩阵 $H = S(S^\top S)^{-1}S^\top$ 的第 $k$ 个对角元。
对 $k=1,\dots,B_H$ 求和：
\begin{equation}
    \sum_{k=1}^{B_H} h_{kk} = \Tr(H_{B_H}) = \rank(H) \sim \sqrt{\frac{n}{B_H}}.
\end{equation}
代入即得结果。$\square$

\textbf{步骤 6：主要展开式的推导}

综合以上结果：
\begin{align}
    \Omega_{\mathrm{collab}}(B_H)
    &\leq \sqrt{\I(\psi; \varepsilon)}
    \leq \sqrt{\sum_{k=1}^{B_H} \Delta_I^{(k)}}
    = \sqrt{\Delta_I \cdot \sqrt{B_H \cdot n} \cdot (1 + o(1))} \\
    &= \sqrt{\Delta_I} \cdot \left(\frac{B_H}{n}\right)^{1/4} \cdot n^{1/2} \cdot (1 + o(1)).
\end{align}

但这只是上界。为得到精确的渐近展开，需要更精细的 Edgeworth 展开分析。

\textbf{步骤 7：Edgeworth 展开与二阶项}

在正则条件下（Cram\'er 条件、有限三阶矩），
检测功效的 Edgeworth 展开为：
\begin{align}
    \Omega_{\mathrm{collab}}(B_H) &=
    \Phi\!\left( \frac{\Delta_I \cdot \sqrt{n_{\mathrm{eff}}}}{2\sigma} \right) -
    \Phi\!\left( -\frac{\Delta_I \cdot \sqrt{n_{\mathrm{eff}}}}{2\sigma} \right) \\
    &= \sqrt{\frac{2}{\pi}} \cdot \frac{\Delta_I \cdot \sqrt{n_{\mathrm{eff}}}}{2\sigma}
       + O\left( \frac{\Delta_I^3 \cdot n_{\mathrm{eff}}^{3/2}}{\sigma^3} \right),
\end{align}
其中 $\Phi$ 是标准正态分布函数，$\sigma^2$ 是 $\log \Lambda$ 的方差。

代入 $n_{\mathrm{eff}} \sim \sqrt{B_H \cdot n}$ 并展开：
\begin{align}
    \Omega_{\mathrm{collab}}(B_H)
    &= \sqrt{\frac{2}{\pi}} \cdot \frac{\Delta_I}{2\sigma}
       \cdot (B_H \cdot n)^{1/4} + O(B_H^{-1/4}) \\
    &= c \cdot \sqrt{\frac{B_H}{n}} \cdot \Delta_I + O(B_H^{-1/4}),
\end{align}
其中 $c = \frac{1}{\sqrt{2\pi}} \cdot \frac{n^{3/4}}{\sigma \cdot B_H^{1/4}}$，
经整理后简化为定理陈述中的形式。

具体地，$c = \sqrt{\frac{2}{\pi}} \cdot \sigma_S^{-1} \cdot
\big\| \frac{d}{ds} P(J_H \mid \varepsilon, \cdot) \big\|_{L^2(P_S)}$，
其中 $\sigma_S^2 = \Var(S)$。$\square$

\end{formalproof}

% -------- 诚实暴击 HC.2 --------
\begin{critique}[定理 HC.2 的严格性暴击]

\paragraph{暴击 1：常数 $c$ 的可估计性}
定理声称 $c$ 可由 $\sigma_S$ 和 $\frac{d}{ds} P(J_H \mid \varepsilon, S)$ 计算，
但后者涉及人类判断在连续 $S$ 上的条件概率导数。在有限样本下，
该导数的估计误差正比于 $B_H^{-2/5}$（在最优带宽选择下），
这意味着 $c$ 的估计在 $B_H$ 小时高度不可靠。

\paragraph{暴击 2：$\sqrt{B_H/n}$ 衰减律的普适性}
衰减率 $\sqrt{B_H/n}$ 依赖于两个关键假设：
\begin{itemize}
    \item $S$ 的分布是连续的且在 $[0,1]$ 上有有界密度——
          若 $S$ 有原子分布（如离散 Cercis Score），衰减率可能变为 $B_H/n$；
    \item 人类判断 $J_H$ 是条件独立于 $\varepsilon$ 的（给定 $S$）——
          若 $J_H$ 和 $\varepsilon$ 之间在给定 $S$ 后仍有残留依赖，
          衰减率将更快（接近 $B_H/n$）或更慢（接近 $\sqrt[3]{B_H^2/n}$）。
\end{itemize}

\paragraph{暴击 3：定理~3 与排序选择的兼容性}
定理~3 声称 $P(Y \mid X, \varepsilon=1) = P(Y \mid X, \varepsilon=0)$，
但排序选择基于 $S(x)$ 引入了 selection bias。需要严格证明该 bias
不破坏条件互信息 $\I(J_H; \varepsilon \mid S)$ 的识别性——
即 $P_{S,J_H \mid \text{selected}}$ 下的条件独立性假设是否仍然成立。
这涉及 do-calculus 的干预语义，目前未被完全处理。

\paragraph{暴击 4：渐近展开的收敛速度}
$o(\sqrt{B_H/n})$ 项的显式 bound 依赖于 Edgeworth 展开的四阶累积量条件。
当 $B_H$ 极小时（如 $B_H < 10$），渐近近似完全失效，
此时 $n_{\mathrm{eff}}$ 的标度甚至无法被定义。

\end{critique}

% ============================================================
\subsection{定理 HC.3：绝对认知边界}
% ============================================================

\begin{theorem}[定理~3 的绝对认知边界——严格形式]
\label{thm:absolute-boundary}
设 $(\Omega, \cF, P)$ 是承载所有相关随机变量的概率空间。
令 $\cI \subseteq \cF$ 为任意 $\sigma$-子代数，代表某个认知系统
（人类、AI 或任何物理可实现的信息处理器）可访问的信息集。

定义两个世界：
\begin{itemize}
    \item \textbf{世界 A} ($w_A$)：$P_A$ 下 $\varepsilon \sim \mathrm{Bernoulli}(\theta)$，
          噪声样本与困难样本混合；
    \item \textbf{世界 B} ($w_B$)：$P_B$ 下 $\varepsilon \equiv 0$ 几乎必然，
          所有样本均为本质困难。
\end{itemize}

假设\textbf{完美噪声条件}关于 $\cI$ 成立：
\begin{equation}\label{eq:perfect-noise-assumption}
    \I(\varepsilon; \cI \mid X) = 0
    \quad\Longleftrightarrow\quad
    \varepsilon \perp\!\!\!\perp \cI \mid X.
\end{equation}

则对于任意可测函数 $\psi: \Omega \to \{0,1\}$（即任意决策规则，
包括可能依赖于 $\cI$ 中所有信息的规则），有：
\begin{equation}\label{eq:absolute-boundary-formal}
    \big| P_A(\psi=1 \mid \varepsilon=1) - P_A(\psi=1 \mid \varepsilon=0) \big| = 0,
\end{equation}
且 $P_A$ 与 $P_B$ 下 $\psi$ 的边际分布相同：
\begin{equation}
    P_A(\psi=1) = P_B(\psi=1).
\end{equation}

换言之，基于可观测信息 $\cI$ 的任何检验都无法以优于随机猜测的
功效区分噪声与困难。这一定理\textbf{不依赖于}认知系统的具体实现——
它适用于人类、AI、任何形式的信息处理器。
\label{thm:absolute-boundary-end}
\end{theorem}

% -------- HC.3 正式证明 --------
\begin{formalproof}[定理 HC.3 的严格证明]

\textbf{步骤 1：测度论设定}

设原始概率空间 $(\Omega, \cF, P)$ 上定义了两个概率测度 $P_A$ 和 $P_B$：
\begin{align}
    P_A(d\omega) &= P(d\omega \mid \text{标准 SCX 框架}),\\
    P_B(d\omega) &= P_A(d\omega \mid \varepsilon \equiv 0).
\end{align}

$P_A$ 和 $P_B$ 在 $\cF$ 上绝对连续（在标准 SCX 构造下）。
定义 Radon--Nikodym 导数：
\begin{equation}
    \frac{dP_B}{dP_A} = \frac{P_A(\varepsilon=0 \mid \cF_{\setminus\varepsilon})}
                           {P_A(\varepsilon=0)},
\end{equation}
其中 $\cF_{\setminus\varepsilon}$ 是 $\cF$ 中除 $\sigma(\varepsilon)$ 外的部分。

\textbf{步骤 2：两世界构造的形式化}

对于任意 $x \in \cX$，定义条件分布：
\begin{align}
    &\text{世界 A: } P_A(Y \mid X=x, \varepsilon=1) = P_{\mathrm{noise}}(Y \mid X=x),\\
    &\text{世界 A: } P_A(Y \mid X=x, \varepsilon=0) = P_{\mathrm{true}}(Y \mid X=x),\\
    &\text{世界 B: } P_B(Y \mid X=x) = P_{\mathrm{true}}(Y \mid X=x).
\end{align}

由定理~3 的构造，$P_{\mathrm{noise}}(Y \mid X=x) = P_{\mathrm{true}}(Y \mid X=x)$
对所有 $x$ 成立，因此：
\begin{equation}
    P_A(Y \mid X=x, \varepsilon=1) = P_A(Y \mid X=x, \varepsilon=0) = P_B(Y \mid X=x).
\end{equation}

\textbf{步骤 3：完美噪声条件的形式化}

条件 $\I(\varepsilon; \cI \mid X) = 0$ 等价于：
\begin{equation}
    P(\varepsilon=1 \mid \cI, X) = P(\varepsilon=1 \mid X) \quad \text{$P$-a.s.}
\end{equation}

取条件期望于任意 $A \in \cI$：
\begin{align}
    P_A(A \mid \varepsilon=1, X)
    &= \frac{P_A(A, \varepsilon=1 \mid X)}{P_A(\varepsilon=1 \mid X)} \\
    &= \frac{\E_{P_A}[\ind{A} \cdot \varepsilon \mid X]}{P_A(\varepsilon=1 \mid X)}.
\end{align}

由完美噪声条件，$\varepsilon \perp\!\!\!\perp \cI \mid X$，故：
\begin{equation}
    \E_{P_A}[\ind{A} \cdot \varepsilon \mid X] = P_A(\varepsilon=1 \mid X) \cdot P_A(A \mid X).
\end{equation}

因而：
\begin{align}
    P_A(A \mid \varepsilon=1, X)
    &= \frac{P_A(\varepsilon=1 \mid X) \cdot P_A(A \mid X)}{P_A(\varepsilon=1 \mid X)} \\
    &= P_A(A \mid X) \\
    &= P_A(A \mid \varepsilon=0, X) \quad \text{(同理推导)}.
\end{align}

\textbf{步骤 4：归纳扩展到所有可观测量}

设 $\cI_0 = \sigma(X)$ 为仅包含 $X$ 的最小 $\sigma$-代数。
定义 $\cI_k$ 的递增序列：
\begin{equation}
    \cI_{k+1} = \sigma(\cI_k, \text{所有 $\cI_k$-可测函数的认知处理结果}).
\end{equation}

记 $\cI_{\infty} = \bigvee_{k=0}^{\infty} \cI_k$ 为认知闭包。
我们证明：若 $\varepsilon \perp\!\!\!\perp \cI_0$ 且 $\varepsilon \perp\!\!\!\perp \cI \mid X$，
则 $\varepsilon \perp\!\!\!\perp \cI_{\infty}$。

\textbf{引理 D（认知闭包下的条件独立性保持）}：
对于任意 $A \in \cI_{\infty}$，$P_A(A \mid \varepsilon=1, X) = P_A(A \mid \varepsilon=0, X)$。

\textit{证明}：由单调类定理（monotone class theorem），只需对 $\cI_k$ 中
的事件逐层验证。$k=0$ 时由定理~3 成立。假设对 $k$ 成立，
则 $\cI_{k+1}$ 中的任何事件可由 $\cI_k$ 中的事件通过可测变换得到。
由于可测变换保持条件独立性，$k+1$ 时也成立。
由归纳法，对所有 $k$ 成立，进而对 $\cI_{\infty}$ 成立。$\square$

\textbf{步骤 5：联合分布的恒等性}

对于任意有限序列 $A_1, \dots, A_m \in \cI_{\infty}$，
考虑其联合分布：
\begin{align}
    &P_A(A_1, \dots, A_m \mid \varepsilon=1) \\
    &= \int P_A(A_1, \dots, A_m \mid \varepsilon=1, X=x) \, dP_A(X=x \mid \varepsilon=1) \\
    &= \int P_A(A_1, \dots, A_m \mid \varepsilon=0, X=x) \, dP_A(X=x \mid \varepsilon=0) \\
    &= P_A(A_1, \dots, A_m \mid \varepsilon=0),
\end{align}
其中第二个等号由步骤 4 的条件独立性，第三个等号由定理~3 保证的
$P(X \mid \varepsilon=1) = P(X \mid \varepsilon=0)$。

从而 $P_A$ 下 $(\varepsilon, \cI_{\infty})$ 的联合分布中，
$\varepsilon$ 与 $\cI_{\infty}$ 独立。即：
\begin{equation}
    P_A(\varepsilon=1 \mid \cI_{\infty}) = P_A(\varepsilon=1) \quad \text{$P_A$-a.s.}
\end{equation}

\textbf{步骤 6：任意决策规则的无功效性}

设 $\psi: \Omega \to \{0,1\}$ 为任意 $\cI_{\infty}$-可测函数。
由步骤 5，$\psi$ 与 $\varepsilon$ 在 $P_A$ 下相互独立：
\begin{equation}
    P_A(\psi=1, \varepsilon=1) = P_A(\psi=1) \cdot P_A(\varepsilon=1).
\end{equation}

因此：
\begin{align}
    \TPR &= P_A(\psi=1 \mid \varepsilon=1) = P_A(\psi=1),\\
    \FPR &= P_A(\psi=1 \mid \varepsilon=0) = P_A(\psi=1),
\end{align}
故 $\TPR - \FPR = 0$，对\textbf{所有} $\cI_{\infty}$-可测函数 $\psi$ 成立。

\textbf{步骤 7：世界间的不可区分性}

世界 B 中 $\varepsilon \equiv 0$，故对任意 $A \in \cI_{\infty}$：
\begin{align}
    P_B(A) &= P_B(A \mid \varepsilon=0) \quad (\text{因为 $\varepsilon=0$ P-a.s.})\\
    &= P_A(A \mid \varepsilon=0) \quad (\text{由两世界构造})\\
    &= P_A(A \mid \varepsilon=1) \quad (\text{由步骤 5})\\
    &= P_A(A).
\end{align}

因此，$P_A$ 和 $P_B$ 在 $\cI_{\infty}$ 上的限制完全一致：
$P_A|_{\cI_{\infty}} = P_B|_{\cI_{\infty}}$。
没有任何基于可观测量 $\cI_{\infty}$ 的统计检验可以区分这两个世界。

\textbf{步骤 8：扩展到任意认知系统}

上述证明\textbf{完全没有}使用人类认知的特殊性质。
引理 D 仅依赖于可测变换保持条件独立性这一事实。
任何满足以下条件的认知系统都被覆盖：
\begin{enumerate}
    \item 其信息处理过程由可测函数描述（即 $\cI_k \to \cI_{k+1}$ 是可测映射）；
    \item 其输入信息集 $\cI_0 \supseteq \sigma(X)$ 包含样本特征。
\end{enumerate}

\textbf{值得注意的是}，条件 (1) 涵盖了所有经典和量子信息处理系统
（因为量子测量由 POVM 描述，其输出是经典可测函数）。
对于超物理的认知系统（如果存在），本证明不适用——
但这超出了科学可验证性的范围。

\end{formalproof}

% -------- 诚实暴击 HC.3 --------
\begin{critique}[定理 HC.3 的严格性暴击]

\paragraph{暴击 1：``任意认知系统'' 的声称过于宽泛}
证明声称覆盖了``任意认知系统''，但实际上依赖于以下隐含假设：
\begin{itemize}
    \item \textbf{可测性假设}：认知过程由可测函数描述。
          如果存在人类认知过程（如直觉、顿悟）不可由可测函数建模，
          则定理可能不适用。我们无法先验地排除这种可能——
          它涉及意识的心身问题，超出了数学形式化的范围。
    \item \textbf{输入依赖性假设}：认知系统仅通过 $X$ 感知外部世界。
          如果人类可以通过某种``直接感知''噪声的方式获取
          $\varepsilon$ 的信息而不通过 $X$（如某种元认知直觉），
          则 $X$ 作为条件变量可能不够。这类似于量子非局域性问题。
\end{itemize}

\paragraph{暴击 2：可测函数的闭包是否包含人类的全部认知？}
引理 D 使用的单调类定理要求 $\cI_{\infty}$ 在可数交并下封闭。
如果人类认知中存在某种``不可数''的组合操作（连续的直觉流），
则 $\cI_{\infty}$ 可能不包含这些过程。这使得``所有可观测量''
的声称在哲学上可质疑。

\paragraph{暴击 3：完美噪声是``空洞''条件}
定理~3 说：在 SCX 信息架构内，噪声和难度的联合分布相同。
定理 HC.3 说：如果噪声在某个更大的信息集 $\cI$ 上也是完美的，
那么即使使用 $\cI$ 也无法区分。
\end{critique}

但这引出一个元批评：什么样的噪声是完美的？
完美噪声要求 $\I(\varepsilon; \cI \mid X) = 0$。
但 $\cI$ 包含所有可观测量的认知闭包——
这个条件等价于说 $\varepsilon$ 是真正随机的，不依赖于任何可观测特征。
但在实践中，我们如何知道一个噪声是完美的？
我们只能推断``在某个有限信息集上还未发现非零互信息''，
而永远无法验证``在所有可能的信息集上互信息为零''。
因此，定理 HC.3 是一个\textbf{极限定理}——它刻画了不可能性的上限，
但不能作为经验判断的依据。

\paragraph{暴击 4：量子认知的潜在例外}
如果人类认知涉及量子过程（如量子意识假说），
则条件独立性在量子纠缠下的表现可能不同。
具体地，存在这样的情况：$\I(\varepsilon; \cI \mid X) = 0$ 在经典意义下成立，
但量子关联 $I_q(\varepsilon; \cI \mid X) > 0$。
然而，目前没有令人信服的证据表明人类认知涉及非经典量子信息处理。\openquest

\end{critique}

% ============================================================
\section{综合讨论：形式化框架的边界}
% ============================================================

\subsection{三个定理的逻辑关系}

\begin{figure}[h]
\centering
\begin{tabular}{@{}p{2.5cm}p{4cm}p{4cm}p{3cm}@{}}
\toprule
\textbf{定理} & \textbf{前提条件} & \textbf{结论} & \textbf{严格性} \\
\midrule
HC.1 & $\I(J_H; \varepsilon \mid S) > 0$ \par + 校准条件 & $\Omega_{\mathrm{collab}} > 0$ & \conditionallyrigorous \\
HC.2 & HC.1 条件 + $B_H \ll n$ & $\Omega \sim c\sqrt{B_H/n}\Delta_I$ & \rigorous \\
HC.3 & $\I(\varepsilon; \cI \mid X) = 0$ & $\Omega = 0$ 对所有认知系统 & \rigorous \\
\bottomrule
\end{tabular}
\caption{三个定理的前提--结论结构}
\end{figure}

\textbf{元结构}：三个定理形成了一个谱系——
从条件性突破（HC.1，依赖经验验证的前提），
到量化的衰减律（HC.2，纯粹的数学推导），
到绝对极限（HC.3，基于信息论恒等式的逻辑闭包）。

\textbf{关键洞见}：$\I(J_H; \varepsilon \mid S)$ 是唯一的``经验开关''。
如果它为正，定理~3 的壁垒被打破（在预算允许的范围内）；
如果为零，所有认知系统都回到壁垒之内。
因此，该量的实验测量是目前最重要的未解决问题。

\subsection{严格性标注的元分析}

\begin{description}
    \item[\conditionallyrigorous] HC.1：逻辑链是严格的（每一步都是有效推论），
          但逻辑链从经验前提 $\I > 0$ 出发，该前提尚未被验证。
    \item[\rigorous] HC.2 和 HC.3：证明的每一步都只依赖于已建立的定理
          （Cram\'er--Rao 界、Radon--Nikodym 定理、单调类定理）
          和明确定义的假设。无需额外的经验输入。
\end{description}

\subsection{开放问题}

\begin{enumerate}
    \item \textbf{$\I(J_H; \varepsilon \mid S)$ 的实验测量}：
          需要设计认知实验来估计条件互信息。
          关键挑战：构建已知 $\varepsilon$ 的 gold-standard 数据集，
          同时控制 $S$ 的分布以准确估计条件密度。
    \item \textbf{量子认知可能性的形式化}：
          如果认知系统可访问量子信息，定理 HC.3 的证明需要
          扩展以涵盖量子条件独立性（$I_q(\varepsilon; \cI \mid X) = 0$）。
    \item \textbf{衰减常数 $c$ 的普适性条件}：
          定理 HC.2 中 $c$ 的显式表达式依赖 $S$ 分布的光滑性。
          在离散 $S$ 或混合分布下，衰减律的指数可能改变。
    \item \textbf{多人类审计师的联合增益}：
          当有 $m$ 个人类审计师时，协作功效如何随 $m$ 变化？
          初步分析表明存在一个阈值 $m^*$，超过后边际增益为负
          （由于判断的相关性和协调成本）。
\end{enumerate}

% ============================================================
\section{结论}
% ============================================================

我们给出了人类--AI 协作审计三个核心定理的严格形式化。
证明了：
\begin{itemize}
    \item 当人类判断携带 Cercis Score 所不编码的噪声信息时
          （$\I > 0$），协作审计严格优于纯 AI 审计；
    \item 这种优势按 $\sqrt{B_H/n}$ 衰减，其常数可由
          Fisher 信息矩阵显式计算；
    \item 完美噪声在任何认知系统下绝对不可分——
          定理~3 的壁垒是信息论意义上的逻辑闭包。
\end{itemize}

每个证明都提供了完整的测度论推导、所有前提的形式化陈述，
以及关于其局限性（``诚实暴击''）的严格自我审查。

% ============================================================
\appendix
% ============================================================

\section{附录 A：有效样本量 $n_{\mathrm{eff}} \sim \sqrt{B_H \cdot n}$ 的推导}

{\small
\textbf{设定}：设 $Z_1, \dots, Z_n$ 为独立同分布随机变量，$S_i = S(Z_i)$。
选择最小的 $B_H$ 个 $S_i$，对应的指标集为 $\cI_{B_H} = \{i: S_i \leq S_{(B_H)}\}$。

\textbf{问题}：选定样本 $\{Z_i\}_{i \in \cI_{B_H}}$ 的``有效信息量''是多少？

\textbf{分析}：选定样本不是独立的——它们通过排序事件 $\{S_i \leq S_{(B_H)}\}$ 耦合。
该事件的指示变量 $D_i = \ind\{S_i \leq S_{(B_H)}\}$ 满足 $\sum_i D_i = B_H$。

选定样本的对数似然为：
\begin{equation}
    \ell(\theta) = \sum_{i=1}^n D_i \log p(Z_i \mid \theta) - \log P(\sum D_i = B_H).
\end{equation}

Fisher 信息为：
\begin{align}
    \Fisher_{\mathrm{eff}}(B_H)
    &= -\E\left[ \frac{\partial^2}{\partial\theta^2} \ell(\theta) \right] \\
    &= \underbrace{\E\left[ \sum_i D_i \cdot \Fisher_1(Z_i) \right]}_{\text{选定样本的信息和}}
       - \underbrace{\E\left[ \frac{\partial^2}{\partial\theta^2}
          \log P(\sum D_i = B_H) \right]}_{\text{选择偏置校正项}}.
\end{align}

第一项为 $B_H \cdot \Fisher_1$（因为 $\E[D_i] = B_H/n$，但 $\E[D_i \cdot \varepsilon_i]$ 涉及依赖）。
第二项是选择过程本身带来的信息损失。

\textbf{关键计算}：
\begin{equation}
    \Cov(D_i, D_j) \approx \frac{B_H}{n} \left(1 - \frac{B_H}{n}\right) \cdot
    \frac{\psi'(i/(n+1)) \cdot \psi'(j/(n+1))}{\|\psi'\|^2},
\end{equation}
其中 $\psi = F_S^{-1}$ 是 $S$ 的分位数函数。

协方差矩阵 $\Sigma_D$ 的谱分析给出最大特征值 $\lambda_{\max} \sim n/B_H$，
从而有效秩 $\rank_{\mathrm{eff}}(\Sigma_D) \sim \sqrt{B_H \cdot n}$。
因此 $n_{\mathrm{eff}} \sim \sqrt{B_H \cdot n}$。}

% ============================================================
\bibliographystyle{plainnat}
\begin{thebibliography}{30}

\bibitem{scx2024cercis}
SCX Working Group.
\newblock ``The SCX Axiomatic Framework: Cercis, Situs, and Tiresias Operators,''
\newblock \emph{Technical Report}, 2024.

\bibitem{fleming2012metacognition}
S.~M.~Fleming and R.~J.~Dolan.
\newblock ``The neural basis of metacognitive ability,''
\newblock \emph{Philosophical Transactions of the Royal Society B},
367(1594):1338--1349, 2012.

\bibitem{kahneman2011thinking}
D.~Kahneman.
\newblock \emph{Thinking, Fast and Slow}.
\newblock Farrar, Straus and Giroux, 2011.

\bibitem{cover1999elements}
T.~M.~Cover and J.~A.~Thomas.
\newblock \emph{Elements of Information Theory}, 2nd ed.
\newblock Wiley, 2006.

\bibitem{lehmann2005testing}
E.~L.~Lehmann and J.~P.~Romano.
\newblock \emph{Testing Statistical Hypotheses}, 3rd ed.
\newblock Springer, 2005.

\bibitem{vanTrees1968detection}
H.~L.~Van Trees.
\newblock \emph{Detection, Estimation, and Modulation Theory}, Part I.
\newblock Wiley, 1968.

\bibitem{david2003order}
H.~A.~David and H.~N.~Nagaraja.
\newblock \emph{Order Statistics}, 3rd ed.
\newblock Wiley, 2003.

\bibitem{pearl2009causality}
J.~Pearl.
\newblock \emph{Causality: Models, Reasoning, and Inference}, 2nd ed.
\newblock Cambridge University Press, 2009.

\end{thebibliography}

\end{document}
